%%
%% Automatically generated file from DocOnce source
%% (https://github.com/doconce/doconce/)
%% doconce format latex week9.do.txt --latex_title_layout=beamer --latex_table_format=footnotesize --no_mako
%%
% #ifdef PTEX2TEX_EXPLANATION
%%
%% The file follows the ptex2tex extended LaTeX format, see
%% ptex2tex: https://code.google.com/p/ptex2tex/
%%
%% Run
%%      ptex2tex myfile
%% or
%%      doconce ptex2tex myfile
%%
%% to turn myfile.p.tex into an ordinary LaTeX file myfile.tex.
%% (The ptex2tex program: https://code.google.com/p/ptex2tex)
%% Many preprocess options can be added to ptex2tex or doconce ptex2tex
%%
%%      ptex2tex -DMINTED myfile
%%      doconce ptex2tex myfile envir=minted
%%
%% ptex2tex will typeset code environments according to a global or local
%% .ptex2tex.cfg configure file. doconce ptex2tex will typeset code
%% according to options on the command line (just type doconce ptex2tex to
%% see examples). If doconce ptex2tex has envir=minted, it enables the
%% minted style without needing -DMINTED.
% #endif

% #define PREAMBLE

% #ifdef PREAMBLE
%-------------------- begin preamble ----------------------

\documentclass[%
oneside,                 % oneside: electronic viewing, twoside: printing
final,                   % draft: marks overfull hboxes, figures with paths
10pt]{article}

\listfiles               %  print all files needed to compile this document

\usepackage{relsize,makeidx,color,setspace,amsmath,amsfonts,amssymb}
\usepackage[table]{xcolor}
\usepackage{bm,ltablex,microtype}

\usepackage[pdftex]{graphicx}

\usepackage{ptex2tex}
% #ifdef MINTED
\usepackage{minted}
\usemintedstyle{default}
% #endif

\usepackage[T1]{fontenc}
%\usepackage[latin1]{inputenc}
\usepackage{ucs}
\usepackage[utf8x]{inputenc}

\usepackage{lmodern}         % Latin Modern fonts derived from Computer Modern

% Hyperlinks in PDF:
\definecolor{linkcolor}{rgb}{0,0,0.4}
\usepackage{hyperref}
\hypersetup{
    breaklinks=true,
    colorlinks=true,
    linkcolor=linkcolor,
    urlcolor=linkcolor,
    citecolor=black,
    filecolor=black,
    %filecolor=blue,
    pdfmenubar=true,
    pdftoolbar=true,
    bookmarksdepth=3   % Uncomment (and tweak) for PDF bookmarks with more levels than the TOC
    }
%\hyperbaseurl{}   % hyperlinks are relative to this root

\setcounter{tocdepth}{2}  % levels in table of contents

% --- fancyhdr package for fancy headers ---
\usepackage{fancyhdr}
\fancyhf{} % sets both header and footer to nothing
\renewcommand{\headrulewidth}{0pt}
\fancyfoot[LE,RO]{\thepage}
% Ensure copyright on titlepage (article style) and chapter pages (book style)
\fancypagestyle{plain}{
  \fancyhf{}
  \fancyfoot[C]{{\footnotesize \copyright\ 1999-2026, Morten Hjorth-Jensen  Email morten.hjorth-jensen@fys.uio.no. Released under CC Attribution-NonCommercial 4.0 license}}
%  \renewcommand{\footrulewidth}{0mm}
  \renewcommand{\headrulewidth}{0mm}
}
% Ensure copyright on titlepages with \thispagestyle{empty}
\fancypagestyle{empty}{
  \fancyhf{}
  \fancyfoot[C]{{\footnotesize \copyright\ 1999-2026, Morten Hjorth-Jensen  Email morten.hjorth-jensen@fys.uio.no. Released under CC Attribution-NonCommercial 4.0 license}}
  \renewcommand{\footrulewidth}{0mm}
  \renewcommand{\headrulewidth}{0mm}
}

\pagestyle{fancy}


\usepackage[framemethod=TikZ]{mdframed}

% --- begin definitions of admonition environments ---

% Admonition style "mdfbox" is an oval colored box based on mdframed
% "notice" admon
\colorlet{mdfbox_notice_background}{gray!5}
\newmdenv[
  skipabove=15pt,
  skipbelow=15pt,
  outerlinewidth=0,
  backgroundcolor=mdfbox_notice_background,
  linecolor=black,
  linewidth=2pt,       % frame thickness
  frametitlebackgroundcolor=mdfbox_notice_background,
  frametitlerule=true,
  frametitlefont=\normalfont\bfseries,
  shadow=false,        % frame shadow?
  shadowsize=11pt,
  leftmargin=0,
  rightmargin=0,
  roundcorner=5,
  needspace=0pt,
]{notice_mdfboxmdframed}

\newenvironment{notice_mdfboxadmon}[1][]{
\begin{notice_mdfboxmdframed}[frametitle=#1]
}
{
\end{notice_mdfboxmdframed}
}

% Admonition style "mdfbox" is an oval colored box based on mdframed
% "summary" admon
\colorlet{mdfbox_summary_background}{gray!5}
\newmdenv[
  skipabove=15pt,
  skipbelow=15pt,
  outerlinewidth=0,
  backgroundcolor=mdfbox_summary_background,
  linecolor=black,
  linewidth=2pt,       % frame thickness
  frametitlebackgroundcolor=mdfbox_summary_background,
  frametitlerule=true,
  frametitlefont=\normalfont\bfseries,
  shadow=false,        % frame shadow?
  shadowsize=11pt,
  leftmargin=0,
  rightmargin=0,
  roundcorner=5,
  needspace=0pt,
]{summary_mdfboxmdframed}

\newenvironment{summary_mdfboxadmon}[1][]{
\begin{summary_mdfboxmdframed}[frametitle=#1]
}
{
\end{summary_mdfboxmdframed}
}

% Admonition style "mdfbox" is an oval colored box based on mdframed
% "warning" admon
\colorlet{mdfbox_warning_background}{gray!5}
\newmdenv[
  skipabove=15pt,
  skipbelow=15pt,
  outerlinewidth=0,
  backgroundcolor=mdfbox_warning_background,
  linecolor=black,
  linewidth=2pt,       % frame thickness
  frametitlebackgroundcolor=mdfbox_warning_background,
  frametitlerule=true,
  frametitlefont=\normalfont\bfseries,
  shadow=false,        % frame shadow?
  shadowsize=11pt,
  leftmargin=0,
  rightmargin=0,
  roundcorner=5,
  needspace=0pt,
]{warning_mdfboxmdframed}

\newenvironment{warning_mdfboxadmon}[1][]{
\begin{warning_mdfboxmdframed}[frametitle=#1]
}
{
\end{warning_mdfboxmdframed}
}

% Admonition style "mdfbox" is an oval colored box based on mdframed
% "question" admon
\colorlet{mdfbox_question_background}{gray!5}
\newmdenv[
  skipabove=15pt,
  skipbelow=15pt,
  outerlinewidth=0,
  backgroundcolor=mdfbox_question_background,
  linecolor=black,
  linewidth=2pt,       % frame thickness
  frametitlebackgroundcolor=mdfbox_question_background,
  frametitlerule=true,
  frametitlefont=\normalfont\bfseries,
  shadow=false,        % frame shadow?
  shadowsize=11pt,
  leftmargin=0,
  rightmargin=0,
  roundcorner=5,
  needspace=0pt,
]{question_mdfboxmdframed}

\newenvironment{question_mdfboxadmon}[1][]{
\begin{question_mdfboxmdframed}[frametitle=#1]
}
{
\end{question_mdfboxmdframed}
}

% Admonition style "mdfbox" is an oval colored box based on mdframed
% "block" admon
\colorlet{mdfbox_block_background}{gray!5}
\newmdenv[
  skipabove=15pt,
  skipbelow=15pt,
  outerlinewidth=0,
  backgroundcolor=mdfbox_block_background,
  linecolor=black,
  linewidth=2pt,       % frame thickness
  frametitlebackgroundcolor=mdfbox_block_background,
  frametitlerule=true,
  frametitlefont=\normalfont\bfseries,
  shadow=false,        % frame shadow?
  shadowsize=11pt,
  leftmargin=0,
  rightmargin=0,
  roundcorner=5,
  needspace=0pt,
]{block_mdfboxmdframed}

\newenvironment{block_mdfboxadmon}[1][]{
\begin{block_mdfboxmdframed}[frametitle=#1]
}
{
\end{block_mdfboxmdframed}
}

% --- end of definitions of admonition environments ---

% prevent orhpans and widows
\clubpenalty = 10000
\widowpenalty = 10000

% --- end of standard preamble for documents ---


% insert custom LaTeX commands...

\raggedbottom
\makeindex

%-------------------- end preamble ----------------------

\begin{document}

% matching end for #ifdef PREAMBLE
% #endif

\newcommand{\exercisesection}[1]{\subsection*{#1}}


% ------------------- main content ----------------------



% ----------------- title -------------------------

\title{Resampling Techniques, Bootstrap and Blocking}

% ----------------- author(s) -------------------------

\author{Morten Hjorth-Jensen  Email morten.hjorth-jensen@fys.uio.no\inst{1}}
\institute{Department of Physics and Center fo Computing in Science Education, University of Oslo, Oslo, Norway\inst{1}}
% ----------------- end author(s) -------------------------

\date{March 20, 2026
% <optional titlepage figure>
% <optional copyright>
}

% !split
\subsection{Overview of week March 16-20, 2025}
\begin{block}{Topics }
\begin{enumerate}
\item Resampling Techniques, Bootstrap and Blocking 

\item Calculation of one-body densities (see whiteboard notes)

\item \href{{https://youtu.be/cnPh5FuGf3o}}{Video of lecture}

\item \href{{https://github.com/CompPhysics/ComputationalPhysics2/blob/gh-pages/doc/HandWrittenNotes/2025/NotesMarch21.pdf}}{Handwritten notes}
\end{enumerate}

\noindent
\end{block}

% !split
\subsection{Why resampling methods ?}
\begin{block}{Statistical analysis }
\begin{itemize}
\item Our simulations can be treated as \emph{computer experiments}. This is particularly the case for Monte Carlo methods

\item The results can be analysed with the same statistical tools as we would use analysing experimental data.

\item As in all experiments, we are looking for expectation values and an estimate of how accurate they are, i.e., possible sources for errors.
\end{itemize}

\noindent
\end{block}    

% !split
\subsection{Statistical analysis}
\begin{block}{}
\begin{itemize}
\item As in other experiments, many numerical  experiments have two classes of errors:
\begin{enumerate}

\item Statistical errors

\item Systematical errors

\end{enumerate}

\noindent
\item Statistical errors can be estimated using standard tools from statistics

\item Systematical errors are method specific and must be treated differently from case to case. 
\end{itemize}

\noindent
\end{block}    

% !split
\subsection{And why do we use such methods?}

As you will see below, due to correlations between various
measurements, we need to evaluate the so-called covariance in order to
establish a proper evaluation of the total variance and the thereby
the standard deviation of a given expectation value.

The covariance however, leads to an evaluation of a double sum over the various stochastic variables. This becomes computationally too expensive to evaluate.
Methods like the Bootstrap, the Jackknife and/or Blocking allow us to circumvent this problem. 

% !split
\subsection{Central limit theorem, reminder from last week}

Last week we derived the central limit theorem with the following assumptions:

\begin{block}{Measurement $i$ }
We assumed that each individual measurement $x_{ij}$ is represented by stochastic variables which independent and identically distributed (iid).
This defined the sample mean of of experiment $i$ with $n$ samples as
\[
\overline{x}_i=\frac{1}{n}\sum_{j} x_{ij}.
\]
and the sample variance
\[
\sigma^2_i=\frac{1}{n}\sum_{j} \left(x_{ij}-\overline{x}_i\right)^2.
\]
\end{block}

% !split
\subsection{Further remarks}

Note that we use $n$ instead of $n-1$ in the definition of
variance. The sample variance and the sample mean are not necessarily equal to
the exact values we would get if we knew the corresponding probability
distribution.

% !split
\subsection{Statistics, more technicalities}
\begin{block}{}
The desired variance
$\mathrm{var}(\overline X_n)$, i.e.~the sample error squared
$\mathrm{err}_X^2$, is given by:
\begin{equation}
\mathrm{err}_X^2 = \mathrm{var}(\overline X_n) = \frac{1}{n^2}
\sum_{ij} \mathrm{cov}(X_i, X_j)
\label{eq:error_exact}
\end{equation}
We see now that in order to calculate the exact error of the sample
with the above expression, we would need the true means
$\mu_{X_i}^{\phantom X}$ of the stochastic variables $X_i$. To
calculate these requires that we know the true multivariate PDF of all
the $X_i$. But this PDF is unknown to us, we have only got the measurements of
one sample. The best we can do is to let the sample itself be an
estimate of the PDF of each of the $X_i$, estimating all properties of
$X_i$ through the measurements of the sample.
\end{block}

% !split
\subsection{Statistics}
\begin{block}{}
Our estimate of $\mu_{X_i}^{\phantom X}$ is then the sample mean $\bar x$
itself, in accordance with the the central limit theorem:
\[
\mu_{X_i}^{\phantom X} = \langle x_i\rangle \approx \frac{1}{n}\sum_{k=1}^n x_k = \bar x
\]
Using $\bar x$ in place of $\mu_{X_i}^{\phantom X}$ we can give an
\emph{estimate} of the covariance in Eq.~(\ref{eq:error_exact})
\[
\mathrm{cov}(X_i, X_j) = \langle (x_i-\langle x_i\rangle)(x_j-\langle x_j\rangle)\rangle
\approx\langle (x_i - \bar x)(x_j - \bar{x})\rangle,
\]
resulting in
\[ 
\frac{1}{n} \sum_{l}^n \left(\frac{1}{n}\sum_{k}^n (x_k -\bar x_n)(x_l - \bar x_n)\right)=\frac{1}{n}\frac{1}{n} \sum_{kl} (x_k -\bar x_n)(x_l - \bar x_n)=\frac{1}{n}\mathrm{cov}(x)
\]
\end{block}

% !split
\subsection{Statistics and sample variance}
\begin{block}{}
By the same procedure we can use the sample variance as an
estimate of the variance of any of the stochastic variables $X_i$
\[
\mathrm{var}(X_i)=\langle x_i - \langle x_i\rangle\rangle \approx \langle x_i - \bar x_n\rangle\nonumber,
\]
which is approximated as 
\begin{equation}
\mathrm{var}(X_i)\approx \frac{1}{n}\sum_{k=1}^n (x_k - \bar x_n)=\mathrm{var}(x)
\label{eq:var_estimate_i_think}
\end{equation}

Now we can calculate an estimate of the error
$\mathrm{err}_X^{\phantom X}$ of the sample mean $\bar x_n$:
\begin{align}
\mathrm{err}_X^2
&=\frac{1}{n^2}\sum_{ij} \mathrm{cov}(X_i, X_j) \nonumber \\
&\approx&\frac{1}{n^2}\sum_{ij}\frac{1}{n}\mathrm{cov}(x) =\frac{1}{n^2}n^2\frac{1}{n}\mathrm{cov}(x)\nonumber\\
&=\frac{1}{n}\mathrm{cov}(x)
\label{eq:error_estimate}
\end{align}
which is nothing but the sample covariance divided by the number of
measurements in the sample.
\end{block}

% !split
\subsection{Statistics, uncorrelated results}
\begin{block}{}

In the special case that the measurements of the sample are
uncorrelated (equivalently the stochastic variables $X_i$ are
uncorrelated) we have that the off-diagonal elements of the covariance
are zero. This gives the following estimate of the sample error:
\[
\mathrm{err}_X^2=\frac{1}{n^2}\sum_{ij} \mathrm{cov}(X_i, X_j) =
\frac{1}{n^2} \sum_i \mathrm{var}(X_i),
\]
resulting in
\begin{equation}
\mathrm{err}_X^2\approx \frac{1}{n^2} \sum_i \mathrm{var}(x)= \frac{1}{n}\mathrm{var}(x)
\label{eq:error_estimate_uncorrel}
\end{equation}
where in the second step we have used Eq.~(\ref{eq:var_estimate_i_think}).
The error of the sample is then just its standard deviation divided by
the square root of the number of measurements the sample contains.
This is a very useful formula which is easy to compute. It acts as a
first approximation to the error, but in numerical experiments, we
cannot overlook the always present correlations.
\end{block}

% !split
\subsection{Statistics, computations}
\begin{block}{}
For computational purposes one usually splits up the estimate of
$\mathrm{err}_X^2$, given by Eq.~(\ref{eq:error_estimate}), into two
parts
\[
\mathrm{err}_X^2 = \frac{1}{n}\mathrm{var}(x) + \frac{1}{n}(\mathrm{cov}(x)-\mathrm{var}(x)),
\]
which equals
\begin{equation}
\frac{1}{n^2}\sum_{k=1}^n (x_k - \bar x_n)^2 +\frac{2}{n^2}\sum_{k<l} (x_k - \bar x_n)(x_l - \bar x_n)
\label{eq:error_estimate_split_up}
\end{equation}
The first term is the same as the error in the uncorrelated case,
Eq.~(\ref{eq:error_estimate_uncorrel}). This means that the second
term accounts for the error correction due to correlation between the
measurements. For uncorrelated measurements this second term is zero.
\end{block}

% !split
\subsection{Statistics, more on computations of errors}
\begin{block}{}
Computationally the uncorrelated first term is much easier to treat
efficiently than the second.
\[
\mathrm{var}(x) = \frac{1}{n}\sum_{k=1}^n (x_k - \bar x_n)^2 =
\left(\frac{1}{n}\sum_{k=1}^n x_k^2\right) - \bar x_n^2
\]
We just accumulate separately the values $x^2$ and $x$ for every
measurement $x$ we receive. The correlation term, though, has to be
calculated at the end of the experiment since we need all the
measurements to calculate the cross terms. Therefore, all measurements
have to be stored throughout the experiment.
\end{block}

% !split
\subsection{Statistics, wrapping up 1}
\begin{block}{}
Let us analyze the problem by splitting up the correlation term into
partial sums of the form:
\[
f_d = \frac{1}{n-d}\sum_{k=1}^{n-d}(x_k - \bar x_n)(x_{k+d} - \bar x_n)
\]
The correlation term of the error can now be rewritten in terms of
$f_d$
\[
\frac{2}{n}\sum_{k<l} (x_k - \bar x_n)(x_l - \bar x_n) =
2\sum_{d=1}^{n-1} f_d
\]
The value of $f_d$ reflects the correlation between measurements
separated by the distance $d$ in the sample samples.  Notice that for
$d=0$, $f$ is just the sample variance, $\mathrm{var}(x)$. If we divide $f_d$
by $\mathrm{var}(x)$, we arrive at the so called \emph{autocorrelation function}
\[
\kappa_d = \frac{f_d}{\mathrm{var}(x)}
\]
which gives us a useful measure of pairwise correlations
starting always at $1$ for $d=0$.
\end{block}

% !split
\subsection{Statistics, final expression}
\begin{block}{}
The sample error (see eq.~(\ref{eq:error_estimate_split_up})) can now be
written in terms of the autocorrelation function:
\begin{align}
\mathrm{err}_X^2 &=
\frac{1}{n}\mathrm{var}(x)+\frac{2}{n}\cdot\mathrm{var}(x)\sum_{d=1}^{n-1}
\frac{f_d}{\mathrm{var}(x)}\nonumber\\ &=&
\left(1+2\sum_{d=1}^{n-1}\kappa_d\right)\frac{1}{n}\mathrm{var}(x)\nonumber\\
&=\frac{\tau}{n}\cdot\mathrm{var}(x)
\label{eq:error_estimate_corr_time}
\end{align}
and we see that $\mathrm{err}_X$ can be expressed in terms the
uncorrelated sample variance times a correction factor $\tau$ which
accounts for the correlation between measurements. We call this
correction factor the \emph{autocorrelation time}:
\begin{equation}
\tau = 1+2\sum_{d=1}^{n-1}\kappa_d
\label{eq:autocorrelation_time}
\end{equation}
\end{block}

% !split
\subsection{Statistics, effective number of correlations}
\begin{block}{}
For a correlation free experiment, $\tau$
equals 1. From the point of view of
eq.~(\ref{eq:error_estimate_corr_time}) we can interpret a sequential
correlation as an effective reduction of the number of measurements by
a factor $\tau$. The effective number of measurements becomes:
\[
n_\mathrm{eff} = \frac{n}{\tau}
\]
To neglect the autocorrelation time $\tau$ will always cause our
simple uncorrelated estimate of $\mathrm{err}_X^2\approx \mathrm{var}(x)/n$ to
be less than the true sample error. The estimate of the error will be
too \emph{good}. On the other hand, the calculation of the full
autocorrelation time poses an efficiency problem if the set of
measurements is very large.
\end{block}

% !split
\subsection{Computing the correlation function}

This code is best seen with the jupyter-notebook












































\bpycod
#!/usr/bin/env python
import numpy as np
import matplotlib.mlab as mlab
import matplotlib.pyplot as plt
import random

# initialize the rng with a seed, simple uniform distribution
random.seed() 
m = 10000
samplefactor = 1.0/m
x = np.zeros(m)   
MeanValue = 0.
VarValue = 0.
for i in range (m):
    value = random.random()
    x[i] = value
    MeanValue += value
    VarValue += value*value

MeanValue *= samplefactor
VarValue *= samplefactor
Variance = VarValue-MeanValue*MeanValue
STDev = np.sqrt(Variance)
print("MeanValue =", MeanValue)
print("Variance =", Variance)
print("Standard deviation =", STDev)

# Computing the autocorrelation function
autocorrelation = np.zeros(m)
darray = np.zeros(m)
for j in range (m):
    sum = 0.0
    darray[j] = j
    for k in range (m-j):
        sum += (x[k]-MeanValue)*(x[k+j]-MeanValue ) 
    autocorrelation[j] = (sum/Variance)*samplefactor
# Visualize results
plt.plot(darray, autocorrelation,'ro')
plt.axis([0,m,-0.2, 1.1])
plt.xlabel(r'$d$')
plt.ylabel(r'$\kappa_d$')
plt.title(r'autocorrelation function for RNG with uniform distribution')
plt.show()

\epycod


% !split
\subsection{Time Auto-correlation Function}
\begin{block}{}
The so-called time-displacement autocorrelation $\phi(t)$ for a quantity $\mathbf{M}$ is given by
\[
\phi(t) = \int dt' \left[\mathbf{M}(t')-\langle \mathbf{M} \rangle\right]\left[\mathbf{M}(t'+t)-\langle \mathbf{M} \rangle\right],
\]
which can be rewritten as 
\[
\phi(t) = \int dt' \left[\mathbf{M}(t')\mathbf{M}(t'+t)-\langle \mathbf{M} \rangle^2\right],
\]
where $\langle \mathbf{M} \rangle$ is the average value and
$\mathbf{M}(t)$ its instantaneous value. We can discretize this function as follows, where we used our
set of computed values $\mathbf{M}(t)$ for a set of discretized times (our Monte Carlo cycles corresponding to moving all electrons?)
\[
\phi(t)  = \frac{1}{t_{\mathrm{max}}-t}\sum_{t'=0}^{t_{\mathrm{max}}-t}\mathbf{M}(t')\mathbf{M}(t'+t)
-\frac{1}{t_{\mathrm{max}}-t}\sum_{t'=0}^{t_{\mathrm{max}}-t}\mathbf{M}(t')\times
\frac{1}{t_{\mathrm{max}}-t}\sum_{t'=0}^{t_{\mathrm{max}}-t}\mathbf{M}(t'+t).
\label{eq:phitf}
\]
\end{block}

% !split
\subsection{Time Auto-correlation Function}
\begin{block}{}
One should be careful with times close to $t_{\mathrm{max}}$, the upper limit of the sums 
becomes small and we end up integrating over a rather small time interval. This means that the statistical
error in $\phi(t)$ due to the random nature of the fluctuations in $\mathbf{M}(t)$ can become large.

One should therefore choose $t \ll t_{\mathrm{max}}$.

Note that the variable $\mathbf{M}$ can be any expectation values of interest.

The time-correlation function gives a measure of the correlation between the various values of the variable 
at a time $t'$ and a time $t'+t$. If we multiply the values of $\mathbf{M}$ at these two different times,
we will get a positive contribution if they are fluctuating in the same direction, or a negative value
if they fluctuate in the opposite direction. If we then integrate over time, or use the discretized version of, the time correlation function $\phi(t)$ should take a non-zero value if the fluctuations are 
correlated, else it should gradually go to zero. For times a long way apart 
the different values of $\mathbf{M}$  are most likely 
uncorrelated and $\phi(t)$ should be zero.
\end{block}

% !split
\subsection{Time Auto-correlation Function}
\begin{block}{}
We can derive the correlation time by observing that our Metropolis algorithm is based on a random
walk in the space of all  possible spin configurations. 
Our probability 
distribution function $\mathbf{\hat{w}}(t)$ after a given number of time steps $t$ could be written as
\[
   \mathbf{\hat{w}}(t) = \mathbf{\hat{W}^t\hat{w}}(0),
\]
with $\mathbf{\hat{w}}(0)$ the distribution at $t=0$ and $\mathbf{\hat{W}}$ representing the 
transition probability matrix. 
We can always expand $\mathbf{\hat{w}}(0)$ in terms of the right eigenvectors of 
$\mathbf{\hat{v}}$ of $\mathbf{\hat{W}}$ as 
\[
    \mathbf{\hat{w}}(0)  = \sum_i\alpha_i\mathbf{\hat{v}}_i,
\]
resulting in 
\[
   \mathbf{\hat{w}}(t) = \mathbf{\hat{W}}^t\mathbf{\hat{w}}(0)=\mathbf{\hat{W}}^t\sum_i\alpha_i\mathbf{\hat{v}}_i=
\sum_i\lambda_i^t\alpha_i\mathbf{\hat{v}}_i,
\]
with $\lambda_i$ the $i^{\mathrm{th}}$ eigenvalue corresponding to  
the eigenvector $\mathbf{\hat{v}}_i$. 
\end{block}

% !split
\subsection{Time Auto-correlation Function}
\begin{block}{}
If we assume that $\lambda_0$ is the largest eigenvector we see that in the limit $t\rightarrow \infty$,
$\mathbf{\hat{w}}(t)$ becomes proportional to the corresponding eigenvector 
$\mathbf{\hat{v}}_0$. This is our steady state or final distribution. 

We can relate this property to an observable like the mean energy.
With the probabilty $\mathbf{\hat{w}}(t)$ (which in our case is the squared trial wave function) we
can write the expectation values as 
\[
 \langle \mathbf{M}(t) \rangle  = \sum_{\mu} \mathbf{\hat{w}}(t)_{\mu}\mathbf{M}_{\mu},
\] 
or as the scalar of a  vector product
 \[
 \langle \mathbf{M}(t) \rangle  = \mathbf{\hat{w}}(t)\mathbf{m},
\] 
with $\mathbf{m}$ being the vector whose elements are the values of $\mathbf{M}_{\mu}$ in its 
various microstates $\mu$.
\end{block}

% !split
\subsection{Time Auto-correlation Function}
\begin{block}{}
We rewrite this relation  as
 \[
 \langle \mathbf{M}(t) \rangle  = \mathbf{\hat{w}}(t)\mathbf{m}=\sum_i\lambda_i^t\alpha_i\mathbf{\hat{v}}_i\mathbf{m}_i.
\] 
If we define $m_i=\mathbf{\hat{v}}_i\mathbf{m}_i$ as the expectation value of
$\mathbf{M}$ in the $i^{\mathrm{th}}$ eigenstate we can rewrite the last equation as
 \[
 \langle \mathbf{M}(t) \rangle  = \sum_i\lambda_i^t\alpha_im_i.
\] 
Since we have that in the limit $t\rightarrow \infty$ the mean value is dominated by the 
the largest eigenvalue $\lambda_0$, we can rewrite the last equation as
 \[
 \langle \mathbf{M}(t) \rangle  = \langle \mathbf{M}(\infty) \rangle+\sum_{i\ne 0}\lambda_i^t\alpha_im_i.
\] 
We define the quantity
\[
   \tau_i=-\frac{1}{log\lambda_i},
\]
and rewrite the last expectation value as
 \[
 \langle \mathbf{M}(t) \rangle  = \langle \mathbf{M}(\infty) \rangle+\sum_{i\ne 0}\alpha_im_ie^{-t/\tau_i}.
\label{eq:finalmeanm}
\] 
\end{block}

% !split
\subsection{Time Auto-correlation Function}
\begin{block}{}

The quantities $\tau_i$ are the correlation times for the system. They control also the auto-correlation function 
discussed above.  The longest correlation time is obviously given by the second largest
eigenvalue $\tau_1$, which normally defines the correlation time discussed above. For large times, this is the 
only correlation time that survives. If higher eigenvalues of the transition matrix are well separated from 
$\lambda_1$ and we simulate long enough,  $\tau_1$ may well define the correlation time. 
In other cases we may not be able to extract a reliable result for $\tau_1$. 
Coming back to the time correlation function $\phi(t)$ we can present a more general definition in terms
of the mean magnetizations $ \langle \mathbf{M}(t) \rangle$. Recalling that the mean value is equal 
to $ \langle \mathbf{M}(\infty) \rangle$ we arrive at the expectation values
\[
\phi(t) =\langle \mathbf{M}(0)-\mathbf{M}(\infty)\rangle \langle \mathbf{M}(t)-\mathbf{M}(\infty)\rangle,
\]
resulting in
\[
\phi(t) =\sum_{i,j\ne 0}m_i\alpha_im_j\alpha_je^{-t/\tau_i},
\]
which is appropriate for all times.
\end{block}

% !split
\subsection{Correlation Time}
\begin{block}{}

If the correlation function decays exponentially
\[ \phi (t) \sim \exp{(-t/\tau)}\]
then the exponential correlation time can be computed as the average
\[   \tau_{\mathrm{exp}}  =  -\langle  \frac{t}{log|\frac{\phi(t)}{\phi(0)}|} \rangle. \]
If the decay is exponential, then
\[  \int_0^{\infty} dt \phi(t)  = \int_0^{\infty} dt \phi(0)\exp{(-t/\tau)}  = \tau \phi(0),\] 
which  suggests another measure of correlation
\[   \tau_{\mathrm{int}} = \sum_k \frac{\phi(k)}{\phi(0)}, \]
called the integrated correlation time.
\end{block}

% !split
\subsection{Resampling methods: Blocking}

The blocking method was made popular by \href{{https://aip.scitation.org/doi/10.1063/1.457480}}{Flyvbjerg and Pedersen (1989)}
and has become one of the standard ways to estimate the variance
$\mathrm{var}(\widehat{\theta})$ for exactly one estimator $\widehat{\theta}$, namely
$\widehat{\theta} = \overline{X}$, the mean value. 

Assume $n = 2^d$ for some integer $d>1$ and $X_1,X_2,\cdots, X_n$ is a stationary time series to begin with. 
Moreover, assume that the series is asymptotically uncorrelated. We switch to vector notation by arranging $X_1,X_2,\cdots,X_n$ in an $n$-tuple. Define:
\begin{align*}
\hat{X} = (X_1,X_2,\cdots,X_n).
\end{align*}

% !split
\subsection{Why blocking?}

The strength of the blocking method is when the number of
observations, $n$ is large. For large $n$, the complexity of dependent
bootstrapping scales poorly, but the blocking method does not,
moreover, it becomes more accurate the larger $n$ is.

% !split
\subsection{Blocking Transformations}
 We now define the blocking transformations. The idea is to take the mean of subsequent
pair of elements from $\bm{X}$ and form a new vector
$\bm{X}_1$. Continuing in the same way by taking the mean of
subsequent pairs of elements of $\bm{X}_1$ we obtain $\bm{X}_2$, and
so on. 
Define $\bm{X}_i$ recursively by:

\begin{align} 
(\bm{X}_0)_k &\equiv (\bm{X})_k \nonumber \\
(\bm{X}_{i+1})_k &\equiv \frac{1}{2}\Big( (\bm{X}_i)_{2k-1} +
(\bm{X}_i)_{2k} \Big) \qquad \text{for all} \qquad 1 \leq i \leq d-1
\end{align} 

% !split
\subsection{Blocking transformations}

The quantity $\bm{X}_k$ is
subject to $k$ \textbf{blocking transformations}.  We now have $d$ vectors
$\bm{X}_0, \bm{X}_1,\cdots,\bm{X}_{d-1}$ containing the subsequent
averages of observations. It turns out that if the components of
$\bm{X}$ is a stationary time series, then the components of
$\bm{X}_i$ is a stationary time series for all $0 \leq i \leq d-1$

We can then compute the autocovariance (or just covariance), the variance, sample mean, and
number of observations for each $i$. 
Let $\gamma_i, \sigma_i^2,
\overline{X}_i$ denote the covariance, variance and average of the
elements of $\bm{X}_i$ and let $n_i$ be the number of elements of
$\bm{X}_i$. It follows by induction that $n_i = n/2^i$. 

% !split
\subsection{Blocking Transformations}

Using the
definition of the blocking transformation and the distributive
property of the covariance, it is clear that since $h =|i-j|$
we can define
\begin{align}
\gamma_{k+1}(h) &= cov\left( ({X}_{k+1})_{i}, ({X}_{k+1})_{j} \right) \nonumber \\
&=  \frac{1}{4}cov\left( ({X}_{k})_{2i-1} + ({X}_{k})_{2i}, ({X}_{k})_{2j-1} + ({X}_{k})_{2j} \right) \nonumber \\
&=  \frac{1}{2}\gamma_{k}(2h) + \frac{1}{2}\gamma_k(2h+1) \hspace{0.1cm} \mathrm{h = 0} \\
&=\frac{1}{4}\gamma_k(2h-1) + \frac{1}{2}\gamma_k(2h) + \frac{1}{4}\gamma_k(2h+1) \quad \mathrm{else}
\end{align}

The quantity $\hat{X}$ is asymptotically uncorrelated by assumption, $\hat{X}_k$ is also asymptotic uncorrelated. Let's turn our attention to the variance of the sample
mean $\mathrm{var}(\overline{X})$. 

% !split
\subsection{Blocking Transformations, getting there}
We have
\begin{align}
\mathrm{var}(\overline{X}_k) = \frac{\sigma_k^2}{n_k} + \underbrace{\frac{2}{n_k} \sum_{h=1}^{n_k-1}\left( 1 - \frac{h}{n_k} \right)\gamma_k(h)}_{\equiv e_k} = \frac{\sigma^2_k}{n_k} + e_k \quad \text{if} \quad \gamma_k(0) = \sigma_k^2. 
\end{align}
The term $e_k$ is called the \textbf{truncation error}: 
\begin{equation}
e_k = \frac{2}{n_k} \sum_{h=1}^{n_k-1}\left( 1 - \frac{h}{n_k} \right)\gamma_k(h). 
\end{equation}
We can show that $\mathrm{var}(\overline{X}_i) = \mathrm{var}(\overline{X}_j)$ for all $0 \leq i \leq d-1$ and $0 \leq j \leq d-1$. 

% !split
\subsection{Blocking Transformations, final expressions}

We can then wrap up
\begin{align}
n_{j+1} \overline{X}_{j+1}  &= \sum_{i=1}^{n_{j+1}} (\hat{X}_{j+1})_i =  \frac{1}{2}\sum_{i=1}^{n_{j}/2} (\hat{X}_{j})_{2i-1} + (\hat{X}_{j})_{2i} \nonumber \\
&= \frac{1}{2}\left[ (\hat{X}_j)_1 + (\hat{X}_j)_2 + \cdots + (\hat{X}_j)_{n_j} \right] = \underbrace{\frac{n_j}{2}}_{=n_{j+1}} \overline{X}_j = n_{j+1}\overline{X}_j. 
\end{align}
By repeated use of this equation we get $\mathrm{var}(\overline{X}_i) = \mathrm{var}(\overline{X}_0) = \mathrm{var}(\overline{X})$ for all $0 \leq i \leq d-1$. This has the consequence that
\begin{align}
\mathrm{var}(\overline{X}) = \frac{\sigma_k^2}{n_k} + e_k \qquad \text{for all} \qquad 0 \leq k \leq d-1. \label{eq:convergence}
\end{align}

% !split
\subsection{More on the blocking method}

Flyvbjerg and Petersen demonstrated that the sequence
$\{e_k\}_{k=0}^{d-1}$ is decreasing, and conjecture that the term
$e_k$ can be made as small as we would like by making $k$ (and hence
$d$) sufficiently large. The sequence is decreasing.
It means we can apply blocking transformations until
$e_k$ is sufficiently small, and then estimate $\mathrm{var}(\overline{X})$ by
$\widehat{\sigma}^2_k/n_k$. 

For an elegant solution and proof of the blocking method, see the recent article of \href{{https://journals.aps.org/pre/abstract/10.1103/PhysRevE.98.043304}}{Marius Jonsson (former MSc student of the Computational Physics group)}.

% !split
\subsection{Example code form last week}

























































































































































































































\bpycod
# 2-electron VMC code for 2dim quantum dot with importance sampling
# Using gaussian rng for new positions and Metropolis- Hastings 
# Added energy minimization
from math import exp, sqrt
from random import random, seed, normalvariate
import numpy as np
import matplotlib.pyplot as plt
from mpl_toolkits.mplot3d import Axes3D
from matplotlib import cm
from matplotlib.ticker import LinearLocator, FormatStrFormatter
from scipy.optimize import minimize
import sys
import os

# Where to save data files
PROJECT_ROOT_DIR = "Results"
DATA_ID = "Results/EnergyMin"

if not os.path.exists(PROJECT_ROOT_DIR):
    os.mkdir(PROJECT_ROOT_DIR)

if not os.path.exists(DATA_ID):
    os.makedirs(DATA_ID)

def data_path(dat_id):
    return os.path.join(DATA_ID, dat_id)

outfile = open(data_path("Energies.dat"),'w')


# Trial wave function for the 2-electron quantum dot in two dims
def WaveFunction(r,alpha,beta):
    r1 = r[0,0]**2 + r[0,1]**2
    r2 = r[1,0]**2 + r[1,1]**2
    r12 = sqrt((r[0,0]-r[1,0])**2 + (r[0,1]-r[1,1])**2)
    deno = r12/(1+beta*r12)
    return exp(-0.5*alpha*(r1+r2)+deno)

# Local energy  for the 2-electron quantum dot in two dims, using analytical local energy
def LocalEnergy(r,alpha,beta):
    
    r1 = (r[0,0]**2 + r[0,1]**2)
    r2 = (r[1,0]**2 + r[1,1]**2)
    r12 = sqrt((r[0,0]-r[1,0])**2 + (r[0,1]-r[1,1])**2)
    deno = 1.0/(1+beta*r12)
    deno2 = deno*deno
    return 0.5*(1-alpha*alpha)*(r1 + r2) +2.0*alpha + 1.0/r12+deno2*(alpha*r12-deno2+2*beta*deno-1.0/r12)

# Derivate of wave function ansatz as function of variational parameters
def DerivativeWFansatz(r,alpha,beta):
    
    WfDer  = np.zeros((2), np.double)
    r1 = (r[0,0]**2 + r[0,1]**2)
    r2 = (r[1,0]**2 + r[1,1]**2)
    r12 = sqrt((r[0,0]-r[1,0])**2 + (r[0,1]-r[1,1])**2)
    deno = 1.0/(1+beta*r12)
    deno2 = deno*deno
    WfDer[0] = -0.5*(r1+r2)
    WfDer[1] = -r12*r12*deno2
    return  WfDer

# Setting up the quantum force for the two-electron quantum dot, recall that it is a vector
def QuantumForce(r,alpha,beta):

    qforce = np.zeros((NumberParticles,Dimension), np.double)
    r12 = sqrt((r[0,0]-r[1,0])**2 + (r[0,1]-r[1,1])**2)
    deno = 1.0/(1+beta*r12)
    qforce[0,:] = -2*r[0,:]*alpha*(r[0,:]-r[1,:])*deno*deno/r12
    qforce[1,:] = -2*r[1,:]*alpha*(r[1,:]-r[0,:])*deno*deno/r12
    return qforce
    

# Computing the derivative of the energy and the energy 
def EnergyDerivative(x0):

    
    # Parameters in the Fokker-Planck simulation of the quantum force
    D = 0.5
    TimeStep = 0.05
    # positions
    PositionOld = np.zeros((NumberParticles,Dimension), np.double)
    PositionNew = np.zeros((NumberParticles,Dimension), np.double)
    # Quantum force
    QuantumForceOld = np.zeros((NumberParticles,Dimension), np.double)
    QuantumForceNew = np.zeros((NumberParticles,Dimension), np.double)

    energy = 0.0
    DeltaE = 0.0
    alpha = x0[0]
    beta = x0[1]
    EnergyDer = 0.0
    DeltaPsi = 0.0
    DerivativePsiE = 0.0 
    #Initial position
    for i in range(NumberParticles):
        for j in range(Dimension):
            PositionOld[i,j] = normalvariate(0.0,1.0)*sqrt(TimeStep)
    wfold = WaveFunction(PositionOld,alpha,beta)
    QuantumForceOld = QuantumForce(PositionOld,alpha, beta)

    #Loop over MC MCcycles
    for MCcycle in range(NumberMCcycles):
        #Trial position moving one particle at the time
        for i in range(NumberParticles):
            for j in range(Dimension):
                PositionNew[i,j] = PositionOld[i,j]+normalvariate(0.0,1.0)*sqrt(TimeStep)+\
                                       QuantumForceOld[i,j]*TimeStep*D
            wfnew = WaveFunction(PositionNew,alpha,beta)
            QuantumForceNew = QuantumForce(PositionNew,alpha, beta)
            GreensFunction = 0.0
            for j in range(Dimension):
                GreensFunction += 0.5*(QuantumForceOld[i,j]+QuantumForceNew[i,j])*\
	                              (D*TimeStep*0.5*(QuantumForceOld[i,j]-QuantumForceNew[i,j])-\
                                      PositionNew[i,j]+PositionOld[i,j])
      
            GreensFunction = exp(GreensFunction)
            ProbabilityRatio = GreensFunction*wfnew**2/wfold**2
            #Metropolis-Hastings test to see whether we accept the move
            if random() <= ProbabilityRatio:
                for j in range(Dimension):
                    PositionOld[i,j] = PositionNew[i,j]
                    QuantumForceOld[i,j] = QuantumForceNew[i,j]
                wfold = wfnew
        DeltaE = LocalEnergy(PositionOld,alpha,beta)
        DerPsi = DerivativeWFansatz(PositionOld,alpha,beta)
        DeltaPsi += DerPsi
        energy += DeltaE
        DerivativePsiE += DerPsi*DeltaE
            
    # We calculate mean values
    energy /= NumberMCcycles
    DerivativePsiE /= NumberMCcycles
    DeltaPsi /= NumberMCcycles
    EnergyDer  = 2*(DerivativePsiE-DeltaPsi*energy)
    return EnergyDer


# Computing the expectation value of the local energy 
def Energy(x0):
    # Parameters in the Fokker-Planck simulation of the quantum force
    D = 0.5
    TimeStep = 0.05
    # positions
    PositionOld = np.zeros((NumberParticles,Dimension), np.double)
    PositionNew = np.zeros((NumberParticles,Dimension), np.double)
    # Quantum force
    QuantumForceOld = np.zeros((NumberParticles,Dimension), np.double)
    QuantumForceNew = np.zeros((NumberParticles,Dimension), np.double)

    energy = 0.0
    DeltaE = 0.0
    alpha = x0[0]
    beta = x0[1]
    #Initial position
    for i in range(NumberParticles):
        for j in range(Dimension):
            PositionOld[i,j] = normalvariate(0.0,1.0)*sqrt(TimeStep)
    wfold = WaveFunction(PositionOld,alpha,beta)
    QuantumForceOld = QuantumForce(PositionOld,alpha, beta)

    #Loop over MC MCcycles
    for MCcycle in range(NumberMCcycles):
        #Trial position moving one particle at the time
        for i in range(NumberParticles):
            for j in range(Dimension):
                PositionNew[i,j] = PositionOld[i,j]+normalvariate(0.0,1.0)*sqrt(TimeStep)+\
                                       QuantumForceOld[i,j]*TimeStep*D
            wfnew = WaveFunction(PositionNew,alpha,beta)
            QuantumForceNew = QuantumForce(PositionNew,alpha, beta)
            GreensFunction = 0.0
            for j in range(Dimension):
                GreensFunction += 0.5*(QuantumForceOld[i,j]+QuantumForceNew[i,j])*\
	                              (D*TimeStep*0.5*(QuantumForceOld[i,j]-QuantumForceNew[i,j])-\
                                      PositionNew[i,j]+PositionOld[i,j])
      
            GreensFunction = exp(GreensFunction)
            ProbabilityRatio = GreensFunction*wfnew**2/wfold**2
            #Metropolis-Hastings test to see whether we accept the move
            if random() <= ProbabilityRatio:
                for j in range(Dimension):
                    PositionOld[i,j] = PositionNew[i,j]
                    QuantumForceOld[i,j] = QuantumForceNew[i,j]
                wfold = wfnew
        DeltaE = LocalEnergy(PositionOld,alpha,beta)
        energy += DeltaE
        if Printout: 
           outfile.write('%f\n' %(energy/(MCcycle+1.0)))            
    # We calculate mean values
    energy /= NumberMCcycles
    return energy

#Here starts the main program with variable declarations
NumberParticles = 2
Dimension = 2
# seed for rng generator 
seed()
# Monte Carlo cycles for parameter optimization
Printout = False
NumberMCcycles= 10000
# guess for variational parameters
x0 = np.array([0.9,0.2])
# Using Broydens method to find optimal parameters
res = minimize(Energy, x0, method='BFGS', jac=EnergyDerivative, options={'gtol': 1e-4,'disp': True})
x0 = res.x
# Compute the energy again with the optimal parameters and increased number of Monte Cycles
NumberMCcycles= 2**19
Printout = True
FinalEnergy = Energy(x0)
EResult = np.array([FinalEnergy,FinalEnergy])
outfile.close()
#nice printout with Pandas
import pandas as pd
from pandas import DataFrame
data ={'Optimal Parameters':x0, 'Final Energy':EResult}
frame = pd.DataFrame(data)
print(frame)

\epycod


% !split
\subsection{Resampling analysis}

The next step is then to use the above data sets and perform a
resampling analysis using the blocking method
The blocking code, based on the article of \href{{https://journals.aps.org/pre/abstract/10.1103/PhysRevE.98.043304}}{Marius Jonsson} is given here


























































\bpycod
# Common imports
import os

# Where to save the figures and data files
DATA_ID = "Results/EnergyMin"

def data_path(dat_id):
    return os.path.join(DATA_ID, dat_id)

infile = open(data_path("Energies.dat"),'r')

from numpy import log2, zeros, mean, var, sum, loadtxt, arange, array, cumsum, dot, transpose, diagonal, sqrt
from numpy.linalg import inv

def block(x):
    # preliminaries
    n = len(x)
    d = int(log2(n))
    s, gamma = zeros(d), zeros(d)
    mu = mean(x)

    # estimate the auto-covariance and variances 
    # for each blocking transformation
    for i in arange(0,d):
        n = len(x)
        # estimate autocovariance of x
        gamma[i] = (n)**(-1)*sum( (x[0:(n-1)]-mu)*(x[1:n]-mu) )
        # estimate variance of x
        s[i] = var(x)
        # perform blocking transformation
        x = 0.5*(x[0::2] + x[1::2])
   
    # generate the test observator M_k from the theorem
    M = (cumsum( ((gamma/s)**2*2**arange(1,d+1)[::-1])[::-1] )  )[::-1]

    # we need a list of magic numbers
    q =array([6.634897,9.210340, 11.344867, 13.276704, 15.086272, 16.811894, 18.475307, 20.090235, 21.665994, 23.209251, 24.724970, 26.216967, 27.688250, 29.141238, 30.577914, 31.999927, 33.408664, 34.805306, 36.190869, 37.566235, 38.932173, 40.289360, 41.638398, 42.979820, 44.314105, 45.641683, 46.962942, 48.278236, 49.587884, 50.892181])

    # use magic to determine when we should have stopped blocking
    for k in arange(0,d):
        if(M[k] < q[k]):
            break
    if (k >= d-1):
        print("Warning: Use more data")
    return mu, s[k]/2**(d-k)


x = loadtxt(infile)
(mean, var) = block(x) 
std = sqrt(var)
import pandas as pd
from pandas import DataFrame
data ={'Mean':[mean], 'STDev':[std]}
frame = pd.DataFrame(data,index=['Values'])
print(frame)


\epycod



% ------------------- end of main content ---------------

% #ifdef PREAMBLE
\end{document}
% #endif

